\section*{Rank-two Boolean compatibility and a fully colliding overlap family}

All Fourier entropies are in bits.  Let $U$ be a finite uniform address
cube, and let $g_u:\{-1,1\}^m\to\{-1,1\}$ be balanced Boolean branches.
Write
\[
 B(u,R)=|U|^{-1/2}\widehat g_u(R),\qquad R\ne\varnothing.
\]

\paragraph{Rank-two compatibility lemma.}
If $\operatorname{rank}B=2$, then there are linearly independent balanced
Boolean functions $g,h$ such that every branch belongs to
\[
 \{g,-g,h,-h\}.
\]
Indeed, choose two independent rows, corresponding after Fourier inversion
to $g,h$.  Every branch has the form $f=ag+bh$.  Since $g,h$ are
independent, $r=gh$ is a nonconstant sign, and therefore assumes both signs.
The Booleanity of $f$ says
\[
 |a+b|=|a-b|=1.
\]
Squaring and subtracting gives $ab=0$, while adding gives
$a^2+b^2=1$.  Hence $(a,b)$ is one of $(\pm1,0),(0,\pm1)$.

Equivalently, there are Boolean address functions $\sigma,\tau$ such that
\[
 F(a,x)=\sigma(a)\left(\frac{1+\tau(a)}2g(x)
                         +\frac{1-\tau(a)}2h(x)\right).
\]
This is an exact Boolean restriction absent from a free rank-two matrix.
It does not by itself pay arbitrary $\sigma,\tau$; doing so would reintroduce
the full-FEI controller obstruction.

\paragraph{Separable two-controller corollary.}
If the address is $(c,b)$, $\tau(c,b)=b$ and $\sigma(c,b)=s(c)$, define
\[
 M_{g,h}(b,x)=\frac{g(x)+h(x)}2
             +b\frac{g(x)-h(x)}2.
\]
Then $M_{g,h}$ is Boolean and
\[
 F(c,b,x)=s(c)M_{g,h}(b,x).
\]
The two factors use disjoint physical blocks, so
\[
 H(F)=H(s)+H(M_{g,h}),\qquad
 I(F)=I(s)+I(M_{g,h}).
\]
Moreover
\[
 I(M_{g,h})=\frac{I(g)+I(h)}2+rac{1-\langle g,h\rangle}{2}.
\]
Consequently independent estimates $H(s)\le C_sI(s)$ and
$H(M_{g,h})\le C_MI(M_{g,h})$ imply
\[
 H(F)\le\max\{C_s,C_M\}I(F).
\]

\paragraph{A fully colliding all-parameter family.}
On three data signs put
\[
 g=\frac{x_1+x_2+x_3-x_1x_2x_3}{2},\qquad
 h=\frac{x_1+x_2-x_3+x_1x_2x_3}{2}.
\]
These are balanced Boolean functions, $\langle g,h\rangle=0$, and their
Fourier supports coincide exactly:
\[
 \operatorname{supp}\widehat g=\operatorname{supp}\widehat h
 =\{\{1\},\{2\},\{3\},\{1,2,3\}\}.
\]
Thus this is maximally non-disjoint spectral overlap.  Nevertheless
\[
 M(b,x)=\frac{x_1+x_2+bx_3-bx_1x_2x_3}{2}
\]
has four Fourier atoms of squared mass $1/4$, whence
\[
 H(M)=2,\qquad I(M)=2.
\]

Let $s_t$ be $-1$ only at the all-minus point of the $t$-cube, and $+1$
elsewhere, and define $F_t(c,b,x)=s_t(c)M(b,x)$.  With
$\delta_t=2^{1-t}$,
\[
 I(s_t)=t\delta_t,
\]
and the standard point-defect calculation gives $H(s_t)<4I(s_t)$.
Therefore, for every $t\ge1$,
\[
 H(F_t)=H(s_t)+2<4I(s_t)+8=4I(F_t),
 \qquad I(F_t)=t2^{1-t}+2.
\]

In the literal $2^{t+1}$-branch selector representation, all three data
coordinates occur in every branch.  Hence worst-case read is $2^{t+1}$,
while activation-weighted congestion is exactly one.  The branch Fourier
matrix has rank two, the rank-name cost is one bit, the type controller costs
one physical root, and $s_t$ has an ordered evaluation tree with expected
root count $2-2^{1-t}<2$.  Thus the total expected controller-root cost is
less than three.  The complete spectral label payment is $H(s_t)+H(M)$;
no $k$-fold Cauchy--Schwarz or disjoint-support assumption is used.

The remaining unpaid step is nonseparable rank-two control: arbitrary
correlated $(\sigma,\tau)$ may carry a hard joint address spectrum.  A valid
extension must bound that joint posterior charge or impose a stable
hypothesis excluding arbitrary FEI-hard controllers.
